\documentclass[11pt]{article}

% Page layout and typography
\usepackage[margin=1in]{geometry}
\usepackage{microtype}
\usepackage{parskip}

% Mathematics
\usepackage{amsmath,amssymb,amsthm,mathtools}
\usepackage{bm}

% Figures and links
\usepackage{graphicx}
\graphicspath{{figures/}}
\usepackage{xcolor}
\usepackage[hidelinks]{hyperref}

% Lists
\usepackage{enumitem}
\setlist{nosep}

% Theorem environments
\theoremstyle{definition}
\newtheorem{definition}{Definition}[section]
\newtheorem{example}[definition]{Example}
\newtheorem{remark}[definition]{Remark}

\theoremstyle{plain}
\newtheorem{theorem}[definition]{Theorem}
\newtheorem{lemma}[definition]{Lemma}
\newtheorem{proposition}[definition]{Proposition}
\newtheorem{corollary}[definition]{Corollary}

% Common number systems
\newcommand{\R}{\mathbb{R}}
\newcommand{\C}{\mathbb{C}}
\newcommand{\Q}{\mathbb{Q}}
\newcommand{\Z}{\mathbb{Z}}
\newcommand{\N}{\mathbb{N}}
\newcommand{\D}{\mathbb{D}}
\newcommand{\T}{\mathbb{T}}
\newcommand{\RS}{\mathbb{S}^2}
\newcommand{\per}{\text{per}}

% Common operators
\DeclareMathOperator{\im}{im}
\DeclareMathOperator{\rank}{rank}
\DeclareMathOperator{\Span}{span}
\DeclareMathOperator{\id}{id}


\title{MATH6313 Lecture Notes}
\author{Caelen Gelineau}
\date{\today}

\begin{document}

\maketitle
\tableofcontents
\newpage

\section{Lecture 1 - 25/08/26}
Calculus review.

\section{Lecture 2 - 27/08/26}
\begin{theorem}[Taylor's theorem]
    Suppose $f \in C^n[a,b]$ that $f^{n+1}$ exists on $[a,b]$ and $x_0 \in [a,b]$. $\forall x \in [a,b], \exists \xi (x)$ between $x_0$ and $x$ s.t.
    \[
        f(x) = P_n(x) + R_n(x)
    \]
    where
    \[
        P_n(x)=f(x_0)+f'(x_0)(x-x_0)+\frac{1}{2!}f''(x_0)(x-x_0)^2+\cdots + \frac{1}{n!} f^{(n)}(x_0) (x-x_0)^n
    \]
    \[
        R_n(x) = \frac{f^{(n+1)}(\xi (x))}{(n+1)!}(x-x_0)^{(n+1)}
    \]
    
\end{theorem}

\section{Lecture 9 - 22/09/26}
\subsection{LU Factorization} 
\begin{definition}
    Lower triangular matrix
    \[
        \begin{pmatrix}
            x & 0 \\
            x & x
        \end{pmatrix}
    \]
    
\end{definition}
\begin{definition}
    Upper triangular matrix
    \[
        \begin{pmatrix}
            x & x \\
            0 & x
        \end{pmatrix}
    \]
    
\end{definition}
Want to write the matrix as
\[
    A = LU
\]
\begin{example}
    \[\begin{pmatrix}
        2 & -1 \\
        -1 & 2
    \end{pmatrix}
    =
    \begin{pmatrix}
        2 & -1 \\
        0 & \frac{3}{2}
    \end{pmatrix}
    = 
    \begin{pmatrix}
        1 & 0 \\
        -\frac{1}{2} & 1
    \end{pmatrix}
    \] 
\end{example}
\begin{example}
Define
\[
    A = \begin{pmatrix}
        1 & 2 & 1 \\
        3 & 8 & 1\\
        0 & 4 & 1
    \end{pmatrix}
\]

\[
    L_{21}(-3) = \begin{pmatrix}
        1 & 0 & 0 \\
        -3 & 1 & 0 \\
        0 & 0 & 1
    \end{pmatrix}
\]
\[
    L_{32}(-2) = \begin{pmatrix}
        1 & 0 & 0 \\
        0 & 1 & 0 \\
        0 & -2 & 1 
    \end{pmatrix}
\]
\[
    L_{21}(-3) A = \begin{pmatrix}
        1 & 0 & 0 \\
        -3 & 1 & 0 \\
        0 & 0 & 1
    \end{pmatrix} \begin{pmatrix}
        1 & 2 & 1 \\
        3 & 8 & 1\\
        0 & 4 & 1
    \end{pmatrix}
\]
Repeat, which results in
\[
  L_{32}(-2) L_{21}(-3) A = \dots = U
\]
Now we get
\[
    A = LU
\]
\[
    L^{-1} = L_{32}(-2)L_{21}(-3)
\]
\[
    L = L_{21}^{-1}(-3) L_{32}^{-1}(-2) = \begin{pmatrix}
        1& 0 & 0 \\
        3 & 1 & 0\\
        0 & 2 & 1
    \end{pmatrix}
\]
\end{example}
\subsection{Summary of LU Factorization method}
Facts: 
\begin{enumerate}
    \item $L_{ij}(-c)$ denotes the lower triangular matrix whose only non-zero entries are $1$'s on
    the diagonal, and $-c$ is in the $(i,j)$ position. Then $A \rightarrow L_{ij}(-c)A$ represents the RREF operation \\"row - row $\times c$".
    \item $[L_{ij}(c)]^{-1} = L_{ij}(-c)$
    \item From the previous example, $L_{21}(3)L_{32}(2) = \begin{pmatrix}
        1 & 0 & 0\\
        3 & 1 & 0\\
        0 & 2 & 1
    \end{pmatrix}$
    which gives us the full lower triangular matrix desired for LU factorization, which you can then use to find $U = L^{-1} A$.
\end{enumerate}
\subsection{Solving \texorpdfstring{$Ax = b$}{Ax = b}}
Given $LU = A$, you can solve $Ax =b$ by
\begin{enumerate}
    \item $Lc = b$ (top to bottom)
    \item $Ux = c$ (bottom to top)
\end{enumerate}

\begin{example}
    \[
        LUx = Lc = b
    \]
    \[
    \begin{pmatrix}
        1 & 1 \\
        3 & -4
    \end{pmatrix} = \begin{pmatrix}
        1 & 0 \\
        3 & 1
    \end{pmatrix}\begin{pmatrix}
        1 & 1 \\
        0 & -7
    \end{pmatrix}
    \]
    Using "$Lc = b$", 
    \[
       \begin{pmatrix}
        1 & 0 \\
        3 & 1
    \end{pmatrix}\begin{pmatrix}
        c_1 \\
        c_2
    \end{pmatrix} =
        \begin{pmatrix}
            3\\
            2
        \end{pmatrix}
    \]
    solving from top equation to bottom equation yields
    \[
        \Rightarrow \begin{cases}
            c_1 = 3  \\
            c_2 = -7 
        \end{cases}
    \]
    Now using "$Ux = c$"
    \[
       \begin{pmatrix}
        1 & 1 \\
        0 & -7
    \end{pmatrix} \begin{pmatrix}
        x_1 \\ x_2
    \end{pmatrix} = \begin{pmatrix}
        3 \\ -7
    \end{pmatrix}
    \]
    solve from bottom equation to top equation yields
       \[
        \Rightarrow \begin{cases}
            c_1 = 2  \\
            c_2 = 1
        \end{cases}
    \] 
    
\end{example}

\subsection{Additional pivoting matrix P}
If you have $PA = LU$, and you can solve $Ax = b$ by the following:
\[
    PAx = Pb
\]
and thus $PA = LU$. Thus, you solve
\[
\begin{cases}
    Lc = Pb\\
    Ux = c
\end{cases}
\]

\begin{example}
    \[
    \begin{pmatrix}
        0 & 1 & 0 \\
        0 & 0 & 1 \\
        1 & 0 & 0
    \end{pmatrix}
    \begin{pmatrix}
        2 & 1& 5\\
        4 & 4 & -4 \\
        1 & 3 & 1
    \end{pmatrix}
    =
    \begin{pmatrix}
        1 & 0 & 0\\
        \frac{1}{4} & 1 & 0\\
        \frac{1}{2} & -\frac{1}{2} & 1
    \end{pmatrix}
    \begin{pmatrix}
        4 & 4 & -4 \\
        0 & 2 & 2 \\
        0 & 0 & 8
    \end{pmatrix}
    \]
    Exercise to the reader.
\end{example}

\subsection{Gaussian Elimination vs LU factorization efficiency}
The computational resource for Gaussian is $\mathcal{O}(\frac{2}{3}n^3k)$, while the computational resource for LU factorization is $\mathcal{O}(2n^2k)$, where 
$k$ is the dimension of the matrix.


\subsubsection{Operation count for back substitution (LU method)}
Suppose you have an $n \times n$ upper triangular matrix.
\[
    \begin{pmatrix}
        a_{11} & a_{12} \dots a_{1n} & b_1 \\
        0 & a_{22} \dots a_{2n} & b_2 \\
        0 & 0 \dots a_{nn} & b_n \\
    \end{pmatrix}
\]
Then we solve
\[
    x_1 = \frac{b_1 - a_{12}x_2 - \dots - a_{1n} x_n}{a_{11}}
\]
\[
    \vdots
\]
\[
    x_n = \frac{b_n}{a_{nn}}
\]
So we have for computing $x_1$, $1$ division operation, $n-1$ multiplications, and $n-1$ additions so we have $2n-1$ total operations.
then for $x_2$, you have $2n-3$ operations, all the way until $1+3+5+\dots+(2n-1) = n^2$ operations for $x_n$.
So we have $\mathcal{O}(n^2)$.

\subsubsection{Operation count for Gaussian elimination}
Suppose you have an $n \times n$ matrix.
\[
    \begin{pmatrix}
        a_{11} & a_{12} \dots a_{1n} & b_1 \\
        a_{21} & a_{22} \dots a_{2n} & b_2 \\
        a_{n1} & a_{n2} \dots a_{nn} & b_n \\
    \end{pmatrix}
\]
The function we want to write is:\\
\verb|for j = 1: n-1|\\
\verb|    eliminate column j| $\longrightarrow$ zero below main diagonal\\\\
which is the same as\\
\verb|for j = 1: n-1|\\
\verb|    for i = j+1:n|\\
\verb|        eliminate a(i,j)|\\
\verb|    end|\\
\verb|end|\\

\newpage
explicitly, the algorithm is \\
\verb|for j = 1: n-1|\\
\verb|    for i = j+1:n|\\
\verb|        mult = a(i,j) / a(j,j)|\\
\verb|        for k = j+1 : n|\\
\verb|          a(i,k) = a(i,k)-mult * a(j,k)|\\
\verb|        end|\\
\verb|    end|\\
\verb|end|\\

To "eliminate $a_{21}$",
\[
    \begin{pmatrix}
        a_{11} & a_{12}& \dots& a_{1n} & b_1 \\
        0 & a_{22} - \frac{a_{21}}{a_{11}}a_{12} & \dots & a_{2n}- \frac{a_{21}}{a_{11}}a_{2n}  & b_2 - \frac{a_{21}}{a_{11}}b_{1} \\
        a_{n1} & a_{n2} &\dots & a_{nn} & b_n \\
    \end{pmatrix}
\] 
which is 1 division, $n$ multiplications, and $n$ additions and subtractions. Thus we have a total of $2n+1$ operations. We have the following table for the number of operations
for each matrix entry. The number of operations for each column decreases by 1 as you move left to right since we have one less entry to worry about.
\[
    \begin{bmatrix}
        0\\
        2n+1 & 0\\
        2n+1 & 2(n-1)+1 & 0\\
        \vdots & \vdots & \vdots\\
        2n+1& 2(n-1)+1 & 2(n-2)+1\\
        2n+1& 2(n-1)+1& 2(n-2)+1 & 0\\
    \end{bmatrix}
\]
Then the sum of total operations is
\[
    \sum_{j=1}^{n-1} \sum_{i=1}^{j} [2(j+1)+1] = \sum_{j=1}^{n-1} [2j(j+1)+j] = \mathcal{O}(\frac{2}{3}n^3)
\]
\subsection{\texorpdfstring{$PA = LU$ factorization}{PA = LU factorization}, or Matrix form of Gaussian elimination w/ partial pivoting}
\begin{example}
    Note, the fractions arent actually there, just to track for $L$.
    \[
        A = \begin{pmatrix}
            2 & 1 & 5\\
            4 & 4 & -4\\
            1 & 3 & 1
        \end{pmatrix}
        \longrightarrow
       \begin{pmatrix}
            4 & 4 & -4\\
            2 & 1 & 5\\
            1 & 3 & 1
        \end{pmatrix} 
        \longrightarrow_{r_2 - \frac{1}{2}r_1, \quad r_3 - \frac{1}{4}r_1}
       \begin{pmatrix}
            4 & 4 & -4\\
            \frac{1}{2} & -1 & 7\\
            \frac{1}{4} & 2 & 2
        \end{pmatrix} 
        \longrightarrow        
       \begin{pmatrix}
            4 & 4 & -4\\
            \frac{1}{4} & 2 & 2\\
            \frac{1}{2} & -1 & 7
        \end{pmatrix} 
        \]
        \[
            \longrightarrow_{r_3 - (-\frac{1}{2}r_2)}
            \begin{pmatrix}
            4 & 4 & -4\\
            \frac{1}{4} & 2 & 2\\
            \frac{1}{2} & -\frac{1}{2} & 8
        \end{pmatrix} 
        \]
        
    or equivalently
    \[
        P_1 A, \quad P_1 = 
        \begin{pmatrix}
            0 & 1 & 0\\
            1 & 0 & 0\\
            0 & 0 & 1
        \end{pmatrix}
    \]
    \[
        M_1P_1A, \quad M_1 = 
        \begin{pmatrix}
            1 & 0 & 0\\
            -\frac{1}{2} & 1 & 0\\
            -\frac{1}{4} & 0 & 1
        \end{pmatrix}
    \]
    \[
        P_2M_1P_1A, \quad P_2 =
        \begin{pmatrix}
            1 & 0 & 0\\
            0 & 0 & 1 \\
            0 & 1 & 0
        \end{pmatrix}
    \]
    \[
        M_2P_2M_1P_1A = U, \quad M_2 = 
        \begin{pmatrix}
            1 & 0 & 0\\
            0 & 1 & 0 \\
            0 & \frac{1}{2} & 1
        \end{pmatrix}
    \]
    thus combining we get
    \[
        U = \begin{pmatrix}
            4 & 4 & -4 \\
            0 & 2 & 2 \\
            0 & 0 & 8
        \end{pmatrix}
    \]
    \[
        L = \begin{pmatrix}
            1 & 0 & 0 \\
            \frac{1}{4} & 1 & 0\\
            \frac{1}{2} & -\frac{1}{2} & 1
        \end{pmatrix}
    \]
    \[
    P = \begin{pmatrix}
        0 & 1 & 0\\
        0 & 0 & 1\\
        1 & 0 & 0
    \end{pmatrix}
    \]
    where $P$ tracks the row swapping operations by multiplying the $P_i$ matrices.\\
    We can equivalently write
    \[
        M_2P_2M_1P_1A = U
    \]
    \[
        M_2P_2M_1P_2^{-1}P_2P_1A = U
    \]
    \[
        M_2P_2M_1P_2^{-1}PA = U
    \]
    \[
        L = (P_2M_1P_2^{-1})M_2^{-1}
    \]
Fact: $P_2M_1P_2^{-1}$ has same structure as $M_1$.
    \subsubsection{Summary}
    \begin{definition}
        A permuatation matrix $P$ is an $n \times n$ matrix consisting of all zeroes except for a single 1 in every row and column.
    \end{definition}
    In general you have
    \[
        M_{m-1}P_{m-1}\cdots M_2P_2M_1P_1A = U \iff M'_{m-1}M'_{m-2}\cdots M'_1 P_{m-1}P_{m-2}\cdots P_1A = U
    \]
    where $M'_{m-1} = M_{m-1}$, $M'_{m-2}=P_{m-1}M_{m-2}P_{m-1}^{-1}, \dots, M'_1 = P_{m-1}P_{m-2}\cdots P_2M_1P_2^{-1}\cdots P_{m-1}^{-1}$.
    
\end{example}

\section{Lecture 10 - 24/09/26}
\subsection{Iterative methods for \texorpdfstring{$Ax=b$}{Ax=b}}
Let $x$ be the exact solution and the initial guess is $x_0$. Define $e := x- x_0$.
\begin{definition}
    The residual is defined as $r := b - Ax_0$.
\end{definition}
Notice that
\[
    Ae = A(x-x_0) = Ax - Ax_0 = b - Ax_0 = r
\]
Thus, the error can be computed if the residual is known:
\[
    e = A^{-1}(b-Ax_0)
\]
\[
    x = x_0 + A^{-1}(b-Ax_0)
\]
however, we do not want to compute $A^{-1}$, so we will find something to replace it, giving the following iterative method:
\[
    x_{k+1} = x_k + B(b-Ax_k)
\]
This turns out to be fixed point iteration:
\[
    x_{k+1} = G(x_k), \quad G(x) = x+B(b-Ax)
\]
Write 
\[
    A = L + D + U
\]
where $D$ is a strictly diagonal matrix and $L,U$ as expected. We can then define
$B := D^{-1}$ or $B := (L+D)^{-1}$.
\subsubsection{Alternate derivation}
\[
    Ax = b
\]
\[
    (L+D+U)x=b
\]
\[
    Lx+Dx+Ux=b 
\]
\[
    Lx_k + Dx_{k+1} + Ux_k = b
\]
\[
    x_{k+1} = D^{-1}(b-Lx_k-Ux_k)
\]

\section{Lecture 11 - 29/09/26}
\subsection{Iterative Methods for Ax = b}
Idea: say the exact solution is $x$ and the initial guess is $x_0$. The error is $e = x-x_0$. The residual is $r = b - Ax_0$.
We can relate the residual to the error by
\[
    Ae = A(x-x_0) = Ax-Ax_0 = b-Ax_0 = r
\]
\[
    e = A^{-1}(b-Ax_0)
\]
\[
    x = x_0 + A^{-1}(b-Ax_0)
\]
The iterative method is then
\[
    x_{k+1} = x_k + B(b-Ax_k)
\]
where $B$ replaces the inverse of $A^{-1}$. Examples for $B$ are
\begin{enumerate}
    \item $B = D^{-1}$
    \item $B = (D+L)^{-1}$
\end{enumerate}
where $A = L+D+U$ where $L$ is a lower matrix, $D$ is a diagonal matrix, and $U$ is a strictly upper matrix.
\subsection{Jacobi Method}
We take $B = D^{-1}$.
\begin{align*}
    x_{k+1} &= x_k + D^{-1}(b-Ax_k) \\
    &= x_k + D^{-1}(b-(L+D+U)x_k) \\
    &= D^{-1}(b-(L+U)x_k)
\end{align*}
Which also can be written as
\[
    Dx_{k+1} = b-(L+U)x_k
\]
or
\[
    Dx_{k+1} + (L+U)x_k = b
\]
\begin{example}
    \[
    \begin{pmatrix}
        2 & -1 & 0\\
        -1 & 2 & -1 \\
        0 & -1 & 2 \\
    \end{pmatrix}
    \begin{pmatrix}
        x\\
        y\\
        z
    \end{pmatrix}=
    \begin{pmatrix}
        0\\
        1\\
        0
    \end{pmatrix}
    \]
    \[
       L= \begin{pmatrix}
            0 & 0 & 0\\
            -1 & 0 & 0\\
            0 & -1 & 0 \\
        \end{pmatrix},\quad
        D = \begin{pmatrix}
            2 & 0 & 0\\
            0 & 2 & 0 \\
            0 & 0 & 2 \\
        \end{pmatrix}, \quad
        U = \begin{pmatrix}
            0 & -1 & 0\\
            0 & 0 & -1 \\
            0 & 0 & 0 \\
        \end{pmatrix}
    \]
    Then we have the system
    \[
        \begin{cases}
            2x_{k+1}-y_k+0 = 0 \\
            -x_k+2y_{k+1}-z_k = 1\\
            0-y_k+2z_{k+1} = 0 
        \end{cases}
    \]
    which solves to
    \[
        \begin{cases}
            x_{k+1} = \frac{y_k}{2}\\
            y_{k+1} = \frac{1}{2}(1+x_k+z_k)\\
            z_{k+1} = \frac{y_k }{2}
        \end{cases}
    \]
\end{example}
\begin{definition}
    The $n \times n$ matrix $A = (a_{ij})$ is strictly diagonally dominant if for any $1 \le i \ge m$, we have
    \[
        |a_{ii}| > \sum_{j\ne i }|a_{ij}| 
    \]
\end{definition}
\begin{example}
    \[
    \begin{pmatrix}
        3 & 1 & -1\\
        2 & -5 & 2 \\
        1 & 6 & 8 \\
    \end{pmatrix}
    \]
    is strictly diagonally dominant, while
    \[
        \begin{pmatrix}
            3 & 2 & 6\\
            1 & 8 & 1 \\
            3 & 2 & -2 \\
        \end{pmatrix}
    \]
    is NOT strictly diagonally dominant.
\end{example}
\begin{theorem}
    If $A$ is strictly diagonally dominant, then $A$ is non-singular ($Ax=b$ has a solution).
    Additionally, for every $b$ and initial guess $x_0$, the Jacobi method converges to the unique solution.
    This also applies to the Gauss-Seidel method.
\end{theorem}

\subsection{Gauss-Seidel method}
For this method, we have $B = (L+D)^{-1}$.
\begin{align*}
    x_{k+1} &= x_k + (L+D)^{-1}(b-Ax-k)\\
    &= x_k + (L+D)^{-1}(b-(L+D+U)x_k)\\
    &= (L+D)^{-1}(b-Ux_k)
\end{align*}
which can be written as
\[
    (L+D)x_{k+1}+Ux_k = b
\]
\begin{example}
    \[
    \begin{pmatrix}
        2 & -1 & 0\\
        -1 & 2 & -1 \\
        0 & -1 & 2 \\
    \end{pmatrix}
    \begin{pmatrix}
        x\\
        y\\
        z
    \end{pmatrix}=
    \begin{pmatrix}
        0\\
        1\\
        0
    \end{pmatrix}
    \]
    \[
       L= \begin{pmatrix}
            0 & 0 & 0\\
            -1 & 0 & 0\\
            0 & -1 & 0 \\
        \end{pmatrix},\quad
        D = \begin{pmatrix}
            2 & 0 & 0\\
            0 & 2 & 0 \\
            0 & 0 & 2 \\
        \end{pmatrix}, \quad
        U = \begin{pmatrix}
            0 & -1 & 0\\
            0 & 0 & -1 \\
            0 & 0 & 0 \\
        \end{pmatrix}
    \]
    Then we have the system
    \[
        \begin{cases}
            2x_{k+1}-y_k+0 = 0 \\
            -x_{k+1}+2y_{k+1}-z_k = 1\\
            0-y_{k+1}+2z_{k+1} = 0 
        \end{cases}
    \]
    \[
    L+U =
        \begin{pmatrix}
            2 & 0 & 0\\
            -1 & 2 & 0 \\
            0 & -1 & 2 \\
        \end{pmatrix}
    \]
    which solves to
    \[
        \begin{cases}
            x_{k+1} = \frac{y_k}{2}\\
            y_{k+1} = \frac{(1+x_{k+1}+z_k)}{2}\\
            z_{k+1} = \frac{y_{k+1}}{2}
        \end{cases}
    \]
    The $(L+D)^{-1}$ is exactly forward substitution method.
\end{example}
\newpage
\begin{lemma}
    If $A$ is strictly diagonally dominant, then $\det(A) \ne 0$, i.e. $A$ is non-singular.   
\end{lemma}

\begin{proof}
    If $\det(A)=0$, then $A$ has zero eigenvalues, i.e. $Ax=0$ for some non-zero $x$.
    Then $\sum_{j=1}^{n}a_{ij}x_j = 0$ for all $i = 1,2,\dots,n$. Take the zero eigenvalue and set it equal to $a_{ii}$.
    \[
        |a_{ii} \cdot x_i| = |a_{ii}||x_i|=\sum_{j\ne i }|a_{ij}x_j| \le \max_{j\ne i}|x_j|\sum_{j\ne i }|a_{ij}| <\max_{j\ne i}|x_j|\cdot|a_{ii}|  
    \]
    \[
        \implies |x_i| < \max_{i \ne j}|x_j|, \quad \forall i = 1,2,\cdots, n
    \]
    which is a contradiction with s.d.d.
\end{proof}

\begin{lemma}
    If the $n \times n$ matrix $A$ has spectral radius $\rho (A) < 1$ (all eigenvalues bounded above by 1), where
    \[
        \rho (A) := \max_{1 \le i \le n}\{|\lambda_i| : \lambda \text{ eigenvalue of } A\}
    \]
    then for any $b$ and $x_0$, $x_{k+1} = Ax_k+b$ converges. In particular, $\lim_{k \to \infty} x_k = x_*$ with $x_* = Ax_* + b$ and is unique.
\end{lemma}
\begin{proof}
    Idea:
    \begin{enumerate}
        \item Uniqueness. Let $x_{*_1} = x_{*_2}$ such that $Ax_{*_i} = Ax_{*_i}+b$. Let $z_* = x_{*_1} - x_{*_2}$. 
        Then $z_* = Az_* \implies 1$ is an eigenvalue of A, but $\rho(A) < 1$, so this is a contradiction.
        \item Consider the geometric series $x_* = b+Ab + A^2b+\cdots$.
        \[
            x_* = Ax_* + b \iff (I-A)x_*=b
        \]
        this Neumann series converges due to $\rho(A)<1$, thus a solution exists. Now, $e_k = x_k-x_*$.
        \[
            Ae_k = A(x_k-x_*)=Ax_k+b-(Ax_*+b) = x_{k+1}-x_*=e_{k+1}
        \]
        then
        \[
            \| e_{k+1}\| \le \|Ae_k\| \le \|A\|\|e_k\| \rightarrow 0
        \]
    \end{enumerate}
    
\end{proof}

\begin{proof}[Proof of Theorem 3.4]
    \begin{align*}
        x_{k+1} &= x_k + D^{-1}(b-Ax_k)\\
        &= (I-D^{-1}A)x_k + D^{-1}b
    \end{align*}
    By Lemma 3.7, it suffices to show that $\rho(I-D^{-1}A) < 1$.
    \[
        (I-D^{-1}A)x = \lambda x
    \]
    \[
        (D-A)x=\lambda Dx
    \]
    \[
        (L+U)x = -\lambda Dx
    \]
   \[
    |\lambda a_{ii}x_i| = |\sum_{i \ne j}a_{ij}x_j| \le \max_{j \ne i}|x_j| \sum_{j \ne i}|a_{ij}| < \max_{j \ne i}|x_j|\cdot|a_{ii} | 
   \]
    Let $|x_k| = \max_{i = 1, 2, \dots, n}|x_i| \ne 0$. Then
    \[
        |\lambda a_{kk}x_k| \le \max_{j \ne k} |x_j||a_{kk}| \le |x_k||a_{kk}| \implies |\lambda| < 1 \implies \rho(I-D^{-1}A) < 1.
    \]
    
\end{proof}

\section{Lecture 12 - 01/10/26}
\begin{remark}
    \[
        A = \begin{pmatrix}
            1 & 4 \\
            \frac{1}{8} & 1
        \end{pmatrix}
    \]
    is not strictly diagonally dominant, but $\rho (I - D^{-1}A) < 1 \implies $ convergence. So s.d.d. is a sufficient but not necessary condition. 
\end{remark}

\subsection{Line search method}
We will only deal with symmetric positive definite matrices.
\[
    A^T = A, \quad  \forall x \ne 0, \; x^T A x > 0
\]
\begin{example}
    \[
        \begin{pmatrix}
            2 & -1 \\
            -1 & 2
        \end{pmatrix}
    \]
    is symmetric positive definite.
\end{example}
General method for line search method:\\
\verb|initial guess x_0|\\
\verb|for k = 0, 1, ...|\\
\verb|  choose search direction d_k|\\
\verb|  choose step size alpha_k|\\
\verb|  set x_(k+1) = x_k + \alpha_k * d_k|\\
\verb|end|

\subsection{Steepest descent}
We choose
\[
    d_k = r_k = b-Ax_k
\]
\[
    \alpha_k = \frac{d_k^T r_k}{d_k^TAd_k}
\]

Consider
\[
    f(x) = \frac{1}{2}x^TAx-b^Tx
\]
We want to minimize $f(x)$, so compute the gradient:
\[
    \nabla f(x) = Ax-b = -r
\]
\[
    a_k = \min_\alpha\{f(x_k + \alpha d_k)\}
\]
So to minimize,
\[
    \frac{\partial f}{\partial \alpha} = \nabla f \cdot d_k=[A(x_k+\alpha d_k)-b] \cdot d_k = (-r_k + \alpha A d_k) \cdot d_k = 0 
\]
\[
    (-r_k+\alpha A d_k) \cdot d_k = 0 \implies \alpha_k = \frac{d_k^T r_k}{d_k^TAd_k} = \frac{r^T_k r_k }{r_k^T A r_k}
\]

\begin{remark}
    We can use a better search direction, the conjugate gradient method.
\end{remark}


\subsection{Interpolation}
Given a set of points, how can we find a polynomial that interpolates those points?
\[
    p(x_i ) = y_i, \quad (x_i, y_i) \text{ are the given data points}
\]
\begin{theorem}
    Given data points $(x_1, y_1), \dots, (x_n, y_n)$, there exists a unique polynomial $p(x)$ of degree $n-1$ or less.
    \[
        p(x_i) = y_i, \quad i = 1,2,\dots, n
    \]
    
\end{theorem}

Goal: compute this polynomial $p(x)$.
\begin{example}
    $(0,1), (2,2)$. 
    \[
        p(x) = a_0 + a_1x
    \]
    \[
    \begin{cases}
        a_0 + 0*a_1 = 1\\
        a_0 + 2a_1 = 2
    \end{cases}
    \implies
    \begin{cases}
        a_0 = 1 \\
        a_1 = \frac{1}{2}
    \end{cases}
    \]
    
\end{example}
\subsection{Lagrange Interpolation}
\begin{example}
    $(0,1), (2,2), (3,4)$. Define the following polynomials:
    \[
        L_1(x) = \frac{(x-x_2)(x-x_3)}{(x_1-x_2)(x_1-x_3)}
    \]
    \[
        L_2(x) = \frac{(x-x_1)(x-x_3)}{(x_2-x_1)(x_2-x_3)}
    \]
    \[
        L_3(x) = \frac{(x-x_1)(x-x_2)}{(x_3-x_1)(x_3-x_2)}
    \]
    Then we can compute    
    \[
        L_1(x_1) = 1,\; L_1(x_2) = 0, \; L_1(x_3)=0
    \]
    \[
        L_2(x_1) = 0,\; L_2(x_2) = 1, \; L_3(x_3)=0
    \]
    \[
        L_3(x_1) = 1,\; L_3(x_2) = 0, \; L_3(x_3)=1
    \]    
    Then we can compute the polynomial as
    \[
        p(x) = y_1L_1(x) + y_2L_2(x)+y_3L_3(x)
    \]
    
\end{example}
\begin{definition}
Given $(x_1, y_1), (x_2, y_2), \dots, (x_n, y_n)$, define
\[
    L_1(x) = \frac{(x-x_2)(x-x_3)\cdots(x-x_n)}{(x_1-x_2)(x_1-x_3)\cdots(x_1-x_n)}
\]
\[
    L_k(x) =\frac{(x-x_1)\cdots(x-x_{k-1})(x-x_{k+1})\cdots (x-x_{n})}{(x_k-x_1)\cdots(x_k-x_{k-1})(x_k-x_{k+1})\cdots(x_k-x_n)} = \frac{\prod_{j \ne k } (x-x_j)}{\prod_{j \ne k}(x_k - x_j) }
\]
These $L_k$ are known as a Lagrange basis. then it follows
\[
    p(x) = y_1L_1(x) + \cdots + y_nL_n(x) = \sum_{i=1}^{n}y_iL_i(x) 
\]
\end{definition}
Question: suppose $y_i = f(x_i)$. We consider the interpolation error $f(x)-p(x)$.

\section{Lecture 13 - 08/10/26}
\begin{remark}
    We have $L_k(x_k) = 1, \; \forall k = 1, 2, \dots, n$ and $L_k(x_j) = 0, \; \forall  j\ne k$
\end{remark}

\begin{theorem}
    Assume $p(x)$ is the Lagrange interpolating polynomial that interpolates \\
    $(x_1, y_1), (x_2, y_2), \dots, (x_n, y_n)$. Define $f(x_i) = y_i$. Then the interpolation error is
    \[
        f(x)-p(x) = \frac{(x-x_1)(x-x_2)\cdots (x-x_n )}{n!}f^{(n)}(\xi)
    \]
    where $\xi$ lies between the smallest and largest of $x_1, x_2, \dots, x_n$.
\end{theorem}

\begin{proof}
    If $x = x_i, \; i = 1,2,\dots, n$, then we are done since it vanishes. If $x \ne x_i$, define for any $t$ between the smallest/largest of $x_i$
    \[
        G(t) = [f(t)-p(t)]-\frac{(t-x_1)(t-x_2)\dots(t-x_n)}{(x-x_1)(x-x_2)\dots(x-x_n)}[f(x)-p(x)]
    \]
    notice that $G(x_i) = 0$ and $G(x) = 0$. By MVT, $G(t)$ has $n+1$ distinct zeroes which implies $G'(t)$ has $n$ distinct zeroes, which cascades to $G^{(n)}(t)$ has a single unique zero $\xi$.
    Thus,
    \[
        G^{(n)}(t)=f^{(n)}(t)-p^{(n)}(t) - \frac{n!}{\prod_{i=1}^{n}x-x_i }[f(x)-p(x)].
    \]
    Notice the $p^{(n)}(t)$ term vanishes because $L$ has degree $n-1$. Then we can compute
    \[
        0 = G^{(n)}(\xi)=f^{(n)}(\xi) -\frac{n!}{\prod_{i=1}^{n}x-x_i }[f(x)-p(x)]. 
    \]
    
\end{proof}

\begin{example}
    $y_i = \sin(x_i)$, $x_1 = 0, x_2 = \frac{\pi}{6}, x_3 = \frac{\pi }{3}, x_4 = \frac{\pi }{2}$.
    \[
        \sin(x) - p(x) = \frac{(x-x_1)(x-x_2)(x-x_3)(x-x_4)}{4!}f^{(4)}(\xi)
    \]
\end{example}

\begin{example}
    Interpolate \[
        f(x) = \frac{1}{1+12x^2}
    \]
    at evenly spaced points in $[-1, 1]$. When computing, the error is large at the endpoints. This is known as the Runge phenomenon.
\end{example}

\subsection{Chebyshev nodes and polynomials}
We want a way to reduce this interpolation error and Runge phenomenon.
\begin{theorem}\label{chebnode}
    The choice of points $x_1, x_2, \dots, x_n$ in $[-1, 1]$ that makes the values of the numerator of the interpolation error
    \[
        \max_{-1 \le x \le 1}|(x-x_1)\cdots (x-x_n)|
    \]
    as small as possible is 
    \[
        x_i = \cos\left[\frac{(2i-1)\pi}{2n} \right], \quad i = 1,2,\dots, n
    \]
    They are called the Chebyshev nodes and the minimum value is $\frac{1}{2^{n-1}}$. The minimum is achieved by
    \[
        (x-x_1)(x-x_2)\cdots(x-x_n) = \frac{1}{2^{n-1}}T_n(x)
    \]
    where $T_n(x)$ denotes the degree $n$ Chebyshev polynomial.
\end{theorem}
\begin{remark}
    The Chebyshev nodes distribute in a way such that there are less points in the middle and more points near the endpoints of the interval.
\end{remark}
\begin{example}
    $f(x) = e^x$ with 5 points. $p_4(x)$ is the degree 4 Chebyshev interpolating polynomial.

    \[
        f(x) - p_4(x) = \frac{(x-x_1)\dots(x-x_5)}{5!}f^{(5)}(\xi)
    \]
    then
    \begin{align*}
        x_1 = \cos \frac{\pi }{10}\\
        x_2 = \cos \frac{3 \pi }{10}\\
        x_3 = \cos \frac{5 \pi}{10}\\
        x_4 =\cos \frac{7 \pi}{10}\\
        x_5 = \cos \frac{9 \pi}{10}\\ 
    \end{align*}
    We also know that 
    \[
        |(x-x_1)\cdots(x-x_n)| \le \frac{1}{2^{n-1}} = \frac{1}{2^4}
    \]
    \[
        |f^{(5)}(\xi)| = |e^\xi| \le e
    \]
    
    \[
        |f(x) - p_4(x)| \le \frac{1}{2^4}\cdot \frac{1}{5!}e \approx 0.00142
    \]
    
\end{example}


\begin{definition}
    The $n$th degree Chebyshev polynomial is given by
    \[
        T_n(x) = \cos(n\arccos(x)), \; -1 \le x \le 1
    \]
\end{definition}
\begin{remark}
    We have the following recursive relation:
    \[
        \begin{cases}
                T_0(x) = 1 \\
                T_1(x) = x\\
                T_{n+1}(x) = 2xT_n(x)-T_{n-1}(x) 
        \end{cases}
    \]
\end{remark}
\begin{proof}
    \[
    T_{n+1}= \cos[(n+1)y] = \cos(ny+y) = \cos(ny)\cos(y)-\sin(ny)\sin(y)
    \]
    \[
        T_{n-1} = \cos(ny) \cos y + \sin(ny)\sin(y)
    \]
   \[
    T_{n+1}+T_{n-1} = 2\cos(ny)\cos(y) = 2xT_n(x)
   \]
\end{proof}
\begin{remark} The following facts are true:
    \begin{enumerate}
        \item $T_n(x)$ is a polynomial.
        \item $\deg T_n = n$ with leading coefficient $2^{n-1}$ is monic.
        \item $T_n(1) =1$ and $T_n(-1)=(-1)^n$.
        \item $\max_{-1 \le x \le 1} |T_n(x)| \le 1$.
        \item Roots of $T_n(x) = 0$ are 
            \begin{align*}
            &\cos(n\arccos(x))=0\\
            &n \arccos x = \frac{\pi}{2} (2i-1)\\
            &x_i = \cos \frac{\pi}{2n}(2i-1), \; i = 1,2,\dots, n
            \end{align*}
            where these $x_i$ are the Chebyshev nodes.
        \item $T_n(x)$ alternates between $-1$ and 1 a total of $n+1$ times.
    \end{enumerate}
    
\end{remark}

\section{Lecture 14 - 13/10/26}

\begin{proof}[Proof of Theorem \ref{chebnode}]
    We prove by way of contradiction. Suppose we have the polynomial 
    \[
        p(x) = (x-\hat{x}_1)(x-\hat{x}_2)\cdots(x-\hat{x}_n)
    \]
    and 
    \[
        \max_{-1 \le x \le 1}|p(x)| < \frac{1}{2^{n-1}}
    \]
    We want to construct
    \[
        Q(x) = p(x) - \frac{1}{2^{n-1}}T_n(x)
    \]
    Since $p(x)$ is bounded above and below by $\pm \frac{1}{2^{n-1}}$ and $T_n(x)$ oscillates between $-1$ and 1 by Fact 3,
    $Q(x)$ alternates between negative and positive values $n+1$ times, thus $Q(x)$ has $n$ distinct roots. But since $p(x)$ is monic with degree $n-1$ and $T_n(x)$ has degree $n-1$, $Q(x)$ has at most degree $n-1$.
    Thus, we have a contradiction. Therefore,
    \[
        \max_{-1 \le x \le 1} |(x-\hat{x}_1)(x-\hat{x}_2)\cdots(x-\hat{x}_n)| \ge \frac{1}{2^{n-1}}
    \]
    for any $\hat{x}$.\\
    If $x_i = \cos(\frac{2i-1}{2n}\pi)$, then
    \[
        \max_{-1 \le x \le 1} |(x-\hat{x}_1)(x-\hat{x}_2)\cdots(x-\hat{x}_n)|  = \frac{1}{2^{n-1}}T_n(x)
    \]
    and
    \[
        \max_{-1 \le x \le 1} |(x-\hat{x}_1)(x-\hat{x}_2)\cdots(x-\hat{x}_n)| = \max_{-1 \le x \le 1} |\frac{1}{2^{n-1}}T_n(x)|  = \frac{1}{2^{n-1}} 
    \]
    
\end{proof}

\subsection{Cubic spline interpolation}
\begin{definition}
Given $n$ data points $(x_1, y_1), \dots (x_n, y_n)$ with strictly increasing $x_i$, a cublic spline $S(x)$ through the data points is a set of cubic
polynomials:
\begin{align*}
    S_1(x) &= y_1 + b_1(x-x_1)+c_1(x-x_1)^2+d_1(x-x_1)^3 \text{ on } [x_1, x_2]\\
    \vdots \\
    S_i(x) &= y_i + b_i(x-x_i)+c_i(x-x_i)^2+d_i(x-x_i)^3 \text{ on } [x_i, x_{i+1}]
\end{align*}   
for $i = 1,2,\dots, n-1$.
\end{definition}

\begin{proposition}
    The following properties hold for $S_i(x)$:
\begin{enumerate}
    \item $S_i(x_i) = y_i$ and $S_i(x_{i+1}) = y_{i+1}$
    \item $S'_{i-1}(x_i) = S_i'(x_i)$
    \item $S''_{i-1}(x_i) = S_i''(x_i)$
    \item The natural spline has the additional property of $S_1''(x_1) = 0$ and  $S_{n-1}'' = 0$ called end conditions.
\end{enumerate}    
\end{proposition}

\begin{proof}
    Property 1 direct computation: \[
        y_{i+1} = y_i + b_i(x_{i+1} - x_i) + c_i(x_{i+1}-x_i)^2 + d_i(x_{i+1} - x_i)^3
    \]
    The Property 2 and 3 computation is similar.\\
    Property 4:
    \begin{align*}
        2c_1 &= 0\\
        2c_{n-1} +6 d_{n-1}(x_n - x_{n-1}) = 0 &= 2c_n
    \end{align*}
    Let $\delta_i = x_{i+1}-x_i$ and $\Delta_i = y_{i+1}-y_i$. this gives us a linear system
    \[
        \begin{pmatrix}
            1 & 0 & 0\\
            \delta_1 & 2\delta_1+2\delta_2 & \delta_2 \\
             & \delta_2 & \ddots & \ddots & \ddots\\
            & & \delta_{n-2} & 2\delta_{n-2}+2\delta_{n-1} & \delta_{n-1}\\
            & & & 0 & 1
        \end{pmatrix}
        \begin{pmatrix}
            c_1\\
            c_2\\
            \vdots \\
            c_n
        \end{pmatrix} = 
        \begin{pmatrix}
            0\\
            3(\frac{\Delta_2}{\delta_2} - \frac{\Delta_1}{\delta_1})\\
            \vdots\\
            3(\frac{\Delta_{n-1}}{\delta_{n-1}} - \frac{\Delta_{n-2}}{\delta_{n-2}})\\
            0
        \end{pmatrix}
    \]
\end{proof}

\begin{example}
    Fine the natural spline through $(0,3), (1,-2), and (2,1)$. We have $x_1 = 0, x_2 = 1, x_3 = 2, y_1 = 3, y_2 = -2, y_3 = 1$.
    Then $\delta_1 = 1, \delta_2 = 1, \Delta_1 = -5, \Delta_2 = 3$. Using the tridiagonal system above, we get
    \[
        \begin{pmatrix}
            1 & 0 & 0\\
            1 & 4 & 1 \\
            0 & 0 & 1 \\
        \end{pmatrix}
        \begin{pmatrix}
            c_1\\
            c_2\\
            c_3
        \end{pmatrix}
        =
        \begin{pmatrix}
            0\\24\\0
        \end{pmatrix}
    \]
    which yields
    \[
        \begin{cases}
            c_1 = 0 \\
            c_2 = 6\\
            c_3 = 0 
        \end{cases}
    \]
    Then we solve for $d_i$ and $b_i$ using the values solved for $c_i$. Then you end with
    \[
        \begin{cases}
            S_1(x) = 3-7x+2x^3 \text{ on } [0,1] \\
            S_2(x) = -2-1(x-1)+6(x-1)^2-2(x-1)^3\text{ on } [1,2] 
        \end{cases}
    \]
    
\end{example}

\end{document}